% SOURCE_AUTHORITY: sga5_fr_workpass.tex, lines 3702--3773 (LF-normalized unit).
% SOURCE_SHA256: 791F4EFFC5E02832D5D77ED03518C8156D6F07E4C8238B03545DB93D883FBB28
\section*{6. Apéndice. Fórmula de Lefschetz para haces coherentes}

En este apartado, $S$ designa un esquema noetheriano. Los $S$-esquemas considerados se supondrán separados y de tipo finito. Recordemos que, según Nagata, todo $S$-morfismo entre tales esquemas es compactificable, es decir, se factoriza como una inmersión abierta seguida de un morfismo propio. Para todo $S$-esquema $X$, denotaremos por $D(X)$ la categoría derivada de la categoría de módulos sobre $\cO_X$.

\subsection*{6.1}
Recordemos ([8] Ap.) que, para todo $S$-morfismo $f:X\to Y$, se dispone de un funtor
\[
f^!:D^+(Y)\longrightarrow D^+(X)
\]
(y de isomorfismos de transitividad $(gf)^! = f^!g^!$ para una compuesta $gf$, con condición de cociclo). Si $f$ es propio, $f^!$ es ``adjunto'' por la derecha del funtor $f_*:D^+(X)_{\mathrm{coh}}\to D^+(Y)_{\mathrm{coh}}$ (el índice ``coh'' designa la subcategoría plena formada por los complejos de cohomología coherente); esta adjunción se precisa en un isomorfismo
\begin{equation}
\tag{6.1.1}
f_*\uRHom(E,f^!F)\xrightarrow{\sim}\uRHom(f_*E,F)
\end{equation}
para $E\in\operatorname{ob}D(X)_{\mathrm{coh}}$ y $F\in\operatorname{ob}D^+(Y)$, definido de manera análoga al isomorfismo (SGA~4 XVIII 3.1.9.6). Las consideraciones de ([8] Ap.5), junto con el teorema de dualidad para el espacio proyectivo, muestran que, si $f$ es liso, de dimensión pura $d$, se tiene un isomorfismo canónico
\begin{equation}
\tag{6.1.2}
f^!L\xleftarrow{\sim} f^*L\otimes\Omega^d_{X/Y}[d];
\end{equation}
si, además, $f$ es propio, la flecha
\begin{equation}
\tag{6.1.3}
\operatorname{Tr}_f:f_*(f^*L\otimes\Omega^d_{X/Y}[d])\longrightarrow L
\end{equation}
deducida, mediante 6.1.2, de la flecha de adjunción $f_*f^!\to I$, coincide con el morfismo de traza ([8] VI).

\subsection*{6.2}
A diferencia de lo que sucede en el caso de los coeficientes discretos, aquí no hay una subcategoría razonable de $D$ estable bajo ``las seis operaciones'' $(f^*,f_*,f^!,f_!,\otimes^{\mathbf L},\uRHom)$. La noción de finitud más natural parece ser la de perfección relativa (SGA~6 III 4). Si $X$ es un $S$-esquema y $L\in\operatorname{ob}D(X)$, recordemos que se dice que $L$ es perfecto relativo a $S$ si $L$ tiene dimensión Tor finita relativa a $S$ y cohomología coherente. Denotaremos por $D(X)_{\mathrm{parf}/S}$ la subcategoría plena de $D(X)$ formada por los complejos perfectos relativos a $S$. En general, no es estable ni bajo $\otimes^{\mathbf L}$ ni bajo $\uRHom$. Sin embargo, si $E$ es perfecto (en sentido absoluto) y $F$ es perfecto relativo a $S$, $\uRHom(E,F)$ es perfecto relativo a $S$ (SGA~6 III 4.9). Por otra parte, sean, para $i=1,2$, $X_i$ un $S$-esquema y $L_i\in\operatorname{ob}D^{-}(X_i)$. Pondremos, como en 1.5,
\[
L_1\otimes_S^{\mathbf L}L_2=Lp_1^*L_1\otimes^{\mathbf L}Lp_2^*L_2,
\]
donde $p_i:X_1\times_SX_2\to X_i$ es la proyección. La consideración de este objeto solo es razonable si $X_1$ y $X_2$ son Tor-independientes sobre $S$, es decir (SGA~6 III 1.5), si $\Tor_i^S(\cO_X,\cO_Y)=0$ para $i\ne0$. Supongamos que así ocurre: entonces, si $L_1$ y $L_2$ son perfectos relativos a $S$, también lo es $L_1\otimes_S^{\mathbf L}L_2$ (SGA~6 III 4.7).

Sea $f:X\to Y$ un $S$-morfismo. Si $f$ es propio, se tiene $f_*D(X)_{\mathrm{parf}/S}\subset D(Y)_{\mathrm{parf}/S}$ (SGA~6 III 4.8). Si $f$ tiene dimensión Tor finita, se tiene $f^*D(Y)_{\mathrm{parf}/S}\subset D(X)_{\mathrm{parf}/S}$ (SGA~6 III 4.5), y $f^!D(Y)_{\mathrm{parf}/S}\subset D(X)_{\mathrm{parf}/S}$ (SGA~6 III 4.9).

Para todo $S$-esquema $a:X\to S$, pondremos
\begin{equation}
\tag{6.2.1}
K_X=a^!\cO_S,
\end{equation}
y denotaremos por $D_X$ (o simplemente por $D$) el funtor $\uRHom(-,K_X)$. Si $X$ tiene dimensión Tor finita sobre $S$, $K_X$ es perfecto relativo a $S$ y, para todo $L\in\operatorname{ob}D(X)_{\mathrm{parf}/S}$, se tiene $DL\in\operatorname{ob}D(X)_{\mathrm{parf}/S}$, y la flecha canónica $L\to DDL$ es un isomorfismo (SGA~6 III 4.9).

Sean $f:X\to Y$ un $S$-morfismo, $L\in\operatorname{ob}D(X)$ y $M\in\operatorname{ob}D(Y)$. Si $f^*M$ está definido (por ejemplo, si $M\in\operatorname{ob}D^-(Y)$, o si $f$ tiene dimensión Tor finita), pondremos
\[
\Hom_f^{(1)}(L,M)=\Hom(f^*M,L),\qquad
\Hom_f^{(2)}(L,M)=\Hom(L,f^*M).
\]
Si $f^!M$ está definido (por ejemplo, si $M\in\operatorname{ob}D^+(Y)$, o si $f$ es liso, o finito y de dimensión Tor finita, o una compuesta de tales morfismos, pues en estos casos $f^!$ se prolonga, mediante 6.1.1 y 6.1.2, a un funtor definido sobre toda la categoría $D(Y)$), pondremos
\[
\Hom_f^{(3)}(L,M)=\Hom(L,f^!M),\qquad
\Hom_f^{(4)}(L,M)=\Hom(f^!M,L).
\]
Los elementos de $\Hom_f^{(i)}(L,M)$ se llamarán, como en 1.4, \emph{flechas de tipo $i$ de $L$ a $M$ sobre $f$} (o $f$-flechas de tipo $i$ de $L$ a $M$). Dejamos al lector la tarea de explicitar la composición de las flechas de tipo $i$, para $i$ fijo, y las categorías fibradas y cofibradas correspondientes. Recordemos únicamente las fórmulas de adjunción parcial
\[
\Hom_f^{(1)}(L,M)=\Hom(M,f_*L)\quad(M\in\operatorname{ob}D^{-}(Y),\ L\in\operatorname{ob}D^+(X)),
\]
\[
\Hom_f^{(3)}(L,M)=\Hom(f_*L,M)\quad(f\text{ propio},\ L\in\operatorname{ob}D(X)_{\mathrm{coh}},\ M\in\operatorname{ob}D^+(Y)).
\]

\subsection*{6.3}
En la situación de 1.6, supongamos que $X_1$ y $X_2$ son Tor-independientes sobre $S$, al igual que $Y_1$ y $Y_2$. Entonces, para $L_i\in\operatorname{ob}D(X_i)_{\mathrm{parf}/S}$, la flecha de Künneth 1.6.4 está definida (úsese el hecho de que la dimensión Tor finita relativa se conserva por imagen directa (SGA~6 III 3.7.1)) y es un isomorfismo (para comprobarlo, se pueden suponer afines $S$ y los $Y_i$ y, mediante un cálculo de Čech, reducirse a suponer también afines $X_1$ y $X_2$; la conclusión es entonces una consecuencia inmediata de las hipótesis de Tor-independencia). Por otra parte, para $M_i\in\operatorname{ob}D(Y_i)_{\mathrm{parf}/S}$, con $f_1^!M_1$ o $f_2^!M_2$ en $D^b$, la flecha de Künneth 1.7.3 está definida\footnote{En el caso propio, o cuando el morfismo se factoriza como una inmersión cerrada seguida de un morfismo liso, la definición de la flecha es inmediata; el caso general se reduce al de tal factorización por descenso cohomológico.}, y es un isomorfismo cuando $f_1$ y $f_2$ tienen dimensión Tor finita, como se comprueba fácilmente mediante el isomorfismo de Künneth de tipo 1.6.4 que acabamos de indicar, conjugado con 6.1.2, y mediante 6.1.1 cuando $f_i$ es una inmersión cerrada. En particular, si $X_1$ y $X_2$ tienen dimensión Tor finita sobre $S$ y son Tor-independientes, 1.7.6 es un isomorfismo.

Se define, como en 2.1, $\uRHom(u,v)$ para $(u,v)$ de tipo $(2,1)$ (resp. $(3,4)$, con $f$ propio para evitar los proobjetos debidos a un $f_!$). Bajo hipótesis adecuadas sobre los grados (véase, por ejemplo, ([9] V 2.3)), están definidas las flechas de evaluación (2.2.2) y de producto tensorial (2.2.3, 2.2.4), y las flechas 2.2.3 y 2.2.4 (siempre que estén definidas) son isomorfismos cuando $E_1$ y $E_2$ son perfectos (en sentido absoluto).

En la situación de 2.4, supongamos que los pares $(X_1,X_2)$, $(X_1,Y_2)$, $(Y_1,X_2)$ y $(Y_1,Y_2)$ son Tor-independientes sobre $S$; que $f_1$ es propio; que $E_i,F_i$, $P_i,Q_i$, $\uRHom(E_i,F_i)$ y $\uRHom(P_i,Q_i)$ son perfectos relativos a $S$; y que las flechas de Künneth
\[
f_1^!Q_1\otimes_S^{\mathbf L}Q_2\to(f_1\times_SY_2)^!(Q_1\otimes_S^{\mathbf L}Q_2),\qquad
f_1^!Q_1\otimes_S^{\mathbf L}F_2\to(f_1\times_SY_2)^!(Q_1\otimes_S^{\mathbf L}F_2)
\]
son isomorfismos. Entonces 2.4.1 y el diagrama de 2.5 están definidos, y este último es conmutativo (sin embargo, las flechas horizontales no son necesariamente isomorfismos). La comprobación es la misma que para 2.5.

